\documentclass[11pt]{article}
\usepackage[margin=1.1in]{geometry}
\usepackage{amsmath,amssymb,amsthm}
\usepackage[numbers]{natbib}
\usepackage{hyperref}
\usepackage{booktabs}

\newtheorem{theorem}{Theorem}[section]
\newtheorem{lemma}[theorem]{Lemma}
\newtheorem{conjecture}[theorem]{Conjecture}
\newtheorem{corollary}[theorem]{Corollary}
\theoremstyle{definition}
\newtheorem{observation}[theorem]{Observation}
\newtheorem{remark}[theorem]{Remark}

\newcommand{\dn}[1]{d_{#1}}
\newcommand{\Dnu}{D_{\nu}}
\newcommand{\ord}{\operatorname{ord}}
\newcommand{\den}{\operatorname{den}}
\newcommand{\Z}{\mathbb{Z}}
\newcommand{\Q}{\mathbb{Q}}
\newcommand{\frG}{\mathfrak{G}}
\newcommand{\cF}{\mathcal{F}}

\title{Sharp denominator corrections for Brown--Zudilin cellular\\
linear forms in $\zeta(5)$: measurements, a conjecture,\\ and a lattice
mechanism}
\author{[author list to be determined]\thanks{Draft of 16 July 2026.
Computations, conjectures and proofs in this paper were produced in a
human--AI collaboration; the verification code and complete datasets are
available in the accompanying repository. Correspondence: [address].}}
\date{July 2026}

\begin{document}
\maketitle

\begin{abstract}
Brown and Zudilin [arXiv:2210.03391] constructed effective rational
approximations $I'(\mathbf{a}\cdot n)=Q_n\zeta(5)-P_n$ to $\zeta(5)$ from
cellular integrals on $\overline{\mathcal M}_{0,8}$, together with
experimental integrality statements for the coefficients. By recovering
the exact rational coefficients $P_n,\hat P_n$ at more than $200$
parameter points via certified integer-relation computations, we show
that the published integrality claims are false as stated---their own
printed values violate them---and that the correct statements carry a
correction factor of exactly $12=2^2\cdot 3$ (respectively $2$ for the
$\zeta(3)$-companion), which we conjecture to be sharp and verify in
every measured cell. Unlike the familiar correction factors in the
classical literature (Ap\'ery, Rhin--Viola, Krattenthaler--Rivoal),
which we check are uniformly \emph{removable}, this correction is
\emph{attained}. We propose a mechanism: the coefficients are periods
against a non-primitive integral Betti class produced by the
construction's motivic elimination of $\zeta(2)$, so the correction is a
lattice index bounded by the Bernoulli constant $24$ arising from
$\zeta(2)=-(2\pi i)^2/24$; an elementary elimination-cost lemma, a
conditional theorem, and a heuristic localization in the totally
symmetric case (the factor $12$ appearing as a doubled sine-kernel
Laurent coefficient $\tfrac16$; all symmetric-case laws verified exactly
for $n\le8$) support this. As by-products we determine the
supremum of the Brown--Zudilin ``worthiness'' exponent over the full
parameter cone, verify that the leading factor $2$ in a denominator
conjecture of Krattenthaler--Rivoal is unnecessary for $n\le 5$, and
document an orbit-coherent boundary-prime anomaly. All measured
denominator slack is bounded; no measured direction offers exponential
arithmetic gain.
\end{abstract}

\section{Introduction}

For the linear forms of Ap\'ery-type used in irrationality proofs, the
arithmetic of the rational coefficients is traditionally expressed by
\emph{inclusions}: statements of the shape $D_n\cdot(\text{coefficient})
\in\Z$ for explicit denominators $D_n$ built from $\dn{n}=
\operatorname{lcm}(1,\dots,n)$. These inclusions are proved (or
conjectured) from term-by-term analysis, and the question of their
\emph{sharpness}---whether the true denominators actually attain
$D_n$---is rarely addressed and, to our knowledge, has never been
investigated systematically by direct computation.

This paper does so for the family of cellular integrals
$I(\mathbf{a})$ on $\overline{\mathcal M}_{0,8}$ introduced by Brown and
Zudilin \cite{BZ}, which produce simultaneous rational approximations
$I'=Q\zeta(5)-P$, $I''=Q\zeta(3)-\hat P$ by a motivic elimination of the
parasitic $\zeta(2)$. Our findings:

\begin{enumerate}
\item The integrality statements printed in \cite{BZ} (their
experimental inclusions for $I'$, and the totally symmetric display
$\dn{n}^5 P_n\in\Z$) are \textbf{false as stated}, already against the
paper's own printed values (\S\ref{sec:corrections}).
\item The corrected statements carry a factor of exactly $12$
(resp.\ $2$), and this factor is \textbf{sharp}: it is attained in
$46$ of $164$ systematically sampled cells at the prime $2$ and $30$
cells at the prime $3$, and never exceeded in any of our more than
$200$ certified cells (\S\ref{sec:conjecture}). One isolated,
orbit-coherent, non-persistent exception at a boundary prime $p=7$ is
documented in \S\ref{sec:anomaly}.
\item By contrast, the analogous correction factors appearing in the
classical literature are uniformly \emph{removable}: we verify
$\dn{n}^3 a_n\in\Z$ for Ap\'ery's $\zeta(3)$ forms for $n\le 50$
(sharpening the ubiquitous ``$2\dn{n}^3$''), and we verify that the
leading factor $2$ retained even in the \emph{conjectured} sharp
denominator law of Krattenthaler--Rivoal \cite[eq.~(17.1)]{KR} for
Zudilin's $\zeta(5),\dots,\zeta(11)$ forms is unnecessary for $n\le5$
(\S\ref{sec:crossfamily})---apparently the first direct computation of
those coefficients.
\item We propose a mechanism (\S\ref{sec:mechanism}): the elimination of
$\zeta(2)$ forces the linear form to be the period of a
\emph{non-primitive} integral Betti class, and the identity
$\zeta(2)=-(2\pi i)^2/24$ prices any integral elimination at a lattice
index dividing $24$. An elementary lemma makes this precise; the
measured constants $12$ and $2$ then \emph{measure} $2$-adic indices of
the integral Betti lattice of the underlying rank-$3$ mixed Tate motive.
In the totally symmetric case the excess is localized (heuristically)
in the doubled Laurent coefficient $\tfrac16$ of the reciprocal-sine
kernel; no $p\ge5$ excess is observed through $n=8$, though the support
statement remains conjectural (\S\ref{sec:symmetric}).
\item All measured denominator slack (cells where the true denominator
is \emph{smaller} than predicted) is bounded along every measured
direction, including growth series $n\le8$; consequently the
construction's ``worthiness'' budget receives no hidden arithmetic
subsidy. We also determine $\sup\gamma=0.86597135\ldots$ over the full
$7$-dimensional projective parameter cone---attained exactly at the
authors' published point, whose symmetric coordinates form an arithmetic
progression (\S\ref{sec:gamma}).
\end{enumerate}

Throughout, ``verified'' means: exact integers $Q_n$ from the closed
binomial double sum of \cite{BZ}; values of $I(\mathbf a\cdot n)$
computed to a certified precision (hundreds of digits) from the
$1$-dimensional integral of a product of two ${}_3F_2$'s given in
\cite{BZ}, evaluated by an exact partial-fraction/polylogarithm method;
and $P_n,\hat P_n$ recovered by PSLQ over the basis $\{1,\zeta(2)\}$
with residual certification, precision-stability checks, and exact
agreement with all rational values printed in \cite{BZ}
(\S\ref{sec:method}).

\section{The family and the measurement pipeline}\label{sec:method}

We follow the notation of \cite{BZ}. For admissible integer parameters
$\mathbf a=(a_1,\dots,a_8)$ (the $17$ convergence forms nonnegative; we
require all $28$ forms of the multiset $\mathbf h(\mathbf a)$
nonnegative so that the full symmetry orbit is admissible), the cellular
integral decomposes as
\[
I(\mathbf a) \;=\; 2Q\,(\zeta(5)+2\zeta(3)\zeta(2)) \;-\;
4\hat P\,\zeta(2)\;-\;2P ,\qquad Q\in\Z,\ P,\hat P\in\Q ,
\]
and the object of interest is $I'=Q\zeta(5)-P$. The group
$\frG\cong\Sigma_7$ acts on normalized integrals through the symmetric
parameters $s_0;s_1,\dots,s_7$, permuting the $28$-element multiset
$\mathbf h$, with $I(\mathbf a)/\prod_{i\in\cF}h_i!$ invariant for the
$17$-element index set $\cF$ of \cite{BZ}.

\emph{Exact $Q$.} $Q(\mathbf a\cdot n)$ is computed exactly from the
double binomial sum \cite[eq.~(sumQ)]{BZ}.

\emph{Certified $I$.} $J(\mathbf p;\mathbf q)=I(\mathbf a)$ reduces
\cite{BZ} to a $1$-dimensional integral of a product of two
${}_3F_2$'s. For integer parameters each ${}_3F_2$ coefficient is a
proper rational function of the summation index; its exact partial
fractions over $\Q$ convert every evaluation into finitely many
polylogarithms $\operatorname{Li}_m(z)$, $z\le0$, with rational weights
(analytic continuation built in). Working precision is raised until two
runs agree; tanh--sinh quadrature errors are certified against the
target precision. On validation, this evaluator agrees with the direct
(slow) hypergeometric path to all compared digits ($\ge40$) and is
$\sim57\times$ faster on production cells.

\emph{Recovery.} With $\zeta$-values to matching precision,
$x=(2Q\zeta(5)+4Q\zeta(3)\zeta(2)-I)/2=P+2\hat P\zeta(2)$, and PSLQ on
$(x,1,\zeta(2))$ returns $P,\hat P$ as exact rationals. Certification:
the relation residual must vanish to half the working precision; the
recovery must be stable under a $+80$-digit rerun; and the pipeline must
reproduce the exact printed values of \cite{BZ} in the totally symmetric
case,
\[
Q_1=21,\quad P_1=\tfrac{87}{4},\quad \hat P_1=\tfrac{101}{4},\qquad
Q_2=2989,\quad P_2=\tfrac{1190161}{384},\quad
\hat P_2=\tfrac{344923}{96},
\]
which it does. An incidental end-to-end validation appears in
\S\ref{sec:anomaly}: two independent recoveries at distinct parameter
points agree with the exact factorial-ratio identity
$I'(g\mathbf a)=I'(\mathbf a)\cdot\prod h_i(\mathbf a)!/\prod
h_i(g\mathbf a)!$ through a discrepancy of $7^2$.

\section{The published inclusions are false as stated}
\label{sec:corrections}

Brown--Zudilin state (experimentally) the inclusion
$\dn{m_1n}\cdots \dn{m_5n}\,I'(\mathbf a\cdot n)\in\Z\zeta(5)+\Z$, where
$m_1\ge\cdots\ge m_5$ are the five largest entries of
$\mathbf h(\mathbf a)$, together with a sharpening by the factor
$\Phi_n=\prod_{p>\sqrt{m_1n}}p^{\nu_p}$, and, in the totally symmetric
case, the display $Q_n,\ \dn{n}^2\dn{2n}\hat P_n,\ \dn{n}^5P_n\in\Z$.

Already at $n=2$ in the totally symmetric case, with the paper's own
printed values: $\dn{2}^5P_2=2^5\cdot\frac{1190161}{384}
=\frac{1190161}{12}\notin\Z$ (off by $2^2\cdot3$), and
$\dn{2}^2\dn{4}\hat P_2=48\cdot\frac{344923}{96}
=\frac{344923}{2}\notin\Z$ (off by $2$). The numerators are coprime to
$6$, so the failures are exact. Our recovered values agree with the
printed ones, independently. Analogous failures occur at $3$ of $19$
small interior points in our first systematic sample, always
$2$-adically. The corrected statements are the subject of the next
section.

\section{The sharp correction: data and conjecture}
\label{sec:conjecture}

Extend Brown--Zudilin's sharpening to every prime (legitimate, since the
factorial-ratio identity above is exact $p$-adically):
\[
\Dnu(\mathbf a,n)=\frac{\prod_{i=1}^{5}\dn{m_in}}{\Phi^{*}_n},\qquad
\Phi^{*}_n=\prod_{\text{all }p}p^{\nu_p},\qquad
\nu_p=\max_{g\in\frG}\ \ord_p
\frac{\prod_{i\in\cF}(h_i(\mathbf a)n)!}{\prod_{i\in\cF}(h_i(g\mathbf a)n)!}.
\]

\begin{conjecture}[robust form]\label{conj:main}
For all admissible $\mathbf a$ and $n\ge1$, and every prime $p$,
\[
\ord_p \den(P_n)\ \le\ \ord_p \Dnu(\mathbf a,n)\;+\;
\begin{cases} 2, & p=2,\\ 1, & p=3,\\ 1, & p\ge5,\end{cases}
\]
with the ceilings at $2$ and $3$ attained, the $p\ge5$ term nonzero only
at boundary primes $p=m_1$, and, in the \textbf{strong form}, vanishing
for $n\ge2$.
\end{conjecture}

We emphasize the epistemic status: both forms are falsifiable by a
single future cell, and the boundedness of the $p\ge5$ term---like every
sharpness statement here---rests on finite evidence only. We know of no
structural argument excluding infinitely many boundary-prime
fluctuations; \S\ref{sec:anomaly} documents the one observed orbit and
the failed attempt to make it recur.

\emph{Evidence.} Over $200$ certified cells (systematic random interior
points with entries up to $7$, growth series at selected points with
$n\le8$, and the authors' record point): the $2$-adic excess
distribution over the first $164$ cells is
$\{-6{:}2,\,-5{:}1,\,-4{:}4,\,-3{:}9,\,-2{:}9,\,-1{:}20,\,0{:}33,
\,+1{:}40,\,+2{:}46\}$ --- ceiling $+2$, attained $46$ times, never
exceeded; the $3$-adic distribution is
$\{-4{:}1,\,-3{:}3,\,-2{:}8,\,-1{:}43,\,0{:}133,\,+1{:}30\}$ ---
ceiling $+1$, attained, never exceeded. In the totally symmetric case
the $2$-adic ceiling is attained at every computed $n$
($n=1,\dots,5$ and $8$): $\den(P_n)=4\cdot 2^{\ord_2 \dn{n}^5}\cdot(\text{odd
part})$, and $\den(\hat P_n)=2\cdot2^{\ord_2 \dn{n}^2\dn{2n}}\cdot
(\text{odd part})$ for $n\le5$. The odd parts show only slack (the
$3$-adic prediction over-pays at most points, which is why the $3$-part
of the correction is visible only where the $3$-adic prediction is
tight).

\emph{Growth.} Along every measured series the excess/slack trajectory
is bounded: e.g.\ the largest observed slack point
$\mathbf a=(4,3,4,4,3,2,4,2)$ has $2$-adic slack
$6,5,4,5,3$ bits at $n=1,\dots,5$; points with attained ceiling at
$n=1,2$ can relax into slack at $n=3$. The conjectured bound---not any
pointwise formula---is the stable object.

\section{A boundary-prime anomaly, and orbit coherence}
\label{sec:anomaly}

Exactly one sampled orbit violates the $p\ge5$ clause: at
$\mathbf a=(6,2,6,1,6,4,5,5)$, $n=1$ (where $m_1=7$ is prime and
$\nu_7=1$), $\den(P_1)=2^7\cdot3^4\cdot5^2\cdot7$ exceeds the
$\nu$-sharpened prediction by one factor of $7$. Since the
factorial-ratio identity is exact, this \emph{predicts} that the
maximizing orbit partner $\mathbf a'=(6,2,5,2,4,6,6,5)$ must violate
even the raw ($\nu$-free) law; direct audit confirms
$\den(P_1(\mathbf a'))=2^6\cdot3^4\cdot5^2\cdot7^2$ against a
$d$-product paying a single $7$. Excesses are thus \emph{orbit-coherent}:
the correct invariant is an orbit-level anomaly divisor.

In our samples the effect is rare and non-persistent: a targeted scan
of $40$ fresh points satisfying the same trigger condition
($m_1\in\{5,7,11,13\}$ prime, $\nu_{m_1}\ge1$) produced \emph{no}
further instance, and the anomalous orbit itself relaxes to slack ($-2$
at $p=7$) at $n=2$, at both partners coherently. We stress that this
falsifies the naive ``$\{2,3\}$ only'' clause outright, and that finite
sampling cannot exclude infinitely many such violations across the cone
or along $n$; the robust form of Conjecture~\ref{conj:main} prices them
at one unit at the boundary prime, and even that bound is an empirical
reading, not a theorem. Characterizing the boundary-prime effect exactly
is open, and we regard it as the sharpest known probe of what the
$\frG$-symmetrized inclusion machinery misses.

\section{The mechanism}\label{sec:mechanism}

\begin{lemma}[elimination cost]\label{lem:cost}
Let $e_0,e_2$ span a $2$-dimensional $\Q$-vector space and let
$L$ be a lattice in the dual with basis $\gamma_1,\gamma_2$ pairing as
$\langle\gamma_1,(e_0,e_2)\rangle=(1,\zeta(2))$,
$\langle\gamma_2,(e_0,e_2)\rangle=(0,(2\pi i)^2)$. Every nonzero
$\delta\in L$ with $\langle\delta,e_2\rangle=0$ is an integer multiple
of $\delta_0=24\gamma_1+\gamma_2$; in particular
$\langle\delta,e_0\rangle\in24\Z$.
\end{lemma}

\begin{proof}
$a\zeta(2)+b(2\pi i)^2=0$ and $\zeta(2)/(2\pi i)^2=-\tfrac1{24}$ force
$b=a/24$, so $24\mid a$.
\end{proof}

If the true integral lattice is finer, with $\gamma_2/2^k$ integral, the
cost drops to $24/2^k$. The linear form $I'$ arises exactly by such an
elimination: $\operatorname{Ext}^1_{MT(\Z)_\Q}(\Q(0),\Q(2))=0$
\cite{BrownMTZ} splits the weight-$4$ part rationally, and the identity
$\zeta^{\mathfrak m}(2)=-(\mathbb L^{\mathfrak m})^2/24$ is available
motivically \cite[\S2.2]{BrownP1}. Under two explicitly stated
hypotheses---(H1) the integrand's de Rham coefficients are
$\Dnu$-controlled (the integral form of Brown--Zudilin's own sketched
Barnes analysis), and (H2) a normalization of the integral Betti
lattice---one obtains $24\,\Dnu\,I'\in\Z\zeta(5)+\Z$, with the constant
improving by the lattice index. The measured sharp constants then
\emph{measure} that index: $12=24/2$ for $P$ (equivalently, cost $2$ on
the natural coefficient $2P$), and $2$ for $\hat P$. At the totally
symmetric point the location of the index-$2$ refinement is now
\emph{computed}: all seven iterated residues of the integrand along the
transverse codimension-$2$ strata carrying $\mathrm{gr}^W_4$ equal
$\pm1$ exactly, so the integral de Rham class is primitive and the
factor $2$ is a genuine refinement of the integral Betti lattice
($\gamma_2/2\in L$; index $1$ at $p=3$, consistent with the discharged
$3$-part), localized --- via the $2$-torsion-freeness of the polar
arrangement's integral combinatorics --- in the $B$-relative cellular
structure of the boundary strata. One computation (the $B$-relative
boundary signs) separates this from a fully measurement-independent
derivation. To our knowledge
these are the first experimentally measured integral-lattice invariants
of periods of $\overline{\mathcal M}_{0,8}$; the literature contains the
$\Q$-theory \cite{DG,BrownMTZ,BrownM08} and an empirical
$\{2,3\}$-clearing phenomenon for motivic period expansions
\cite{Schnetz}, but---so far as we can determine---no integral-lattice
synthesis. The support $\{2,3\}$ matches the torsion
$K_3(\Z)\cong\Z/48$ \cite{LS} of the relevant splitting obstruction; we
flag this as a consistency, not a derivation.

\section{The totally symmetric case: concrete derivation and a support
theorem}\label{sec:symmetric}

In the totally symmetric case $\mathbf a=(1,\dots,1)$ the Barnes
representation reduces to the reciprocal-sine kernel with rational part
$r(s)r(t)C(s,t)$, $r(s)=\prod_{i=1}^n(s+i)\big/\prod_{j=n+1}^{2n+1}(s+j)^2$,
whose partial-fraction data is fully explicit:
$B_j=(-1)^n\binom{n+a}{a}\binom{n}{a}^2/n!$ and
$A_j=B_j\,(3H_a-H_{n+a}-2H_{n-a})$ with $a=j-n-1$ (verified exactly for
$n\le6$). Two consequences, proved in the accompanying notes:

\begin{conjecture}[support; verified $n\le8$]\label{thm:support}
In the totally symmetric case every prime dividing $\den(P_n)$ is at
most $2n+1$, and $\ord_p\den(P_n)\le 5\,\ord_p \dn{n}$ for every
$p\ge5$ (inequality only: the corresponding equality fails at
$(n,p)=(8,7)$).
\end{conjecture}

\noindent\emph{Status.} An adversarial audit of our own earlier notes
found that a claimed unconditional proof of this statement silently
invoked the open collapse below, and that even the prime-support bound
is open in general (higher sine-kernel Laurent coefficients
$\zeta(2k)/\pi^{2k}$ carry Bernoulli numerators with large primes,
whose non-contribution is exactly what must be proven). What is proven:
the Barnes specialization and the sine-kernel reduction (complete
analytic proofs), and the classical partial-fraction/harmonic-clearing
integrality inputs. What is open: the general-$n$ ``$\dn{n}^5$-collapse''
of the residue sums (the coupling polynomial of the double kernel admits
no closed-form reduction to the elementary kernels, which is precisely
where an earlier claimed reduction was found to be circular) and the
$2$-adic attainment. All laws --- $12\,\dn{n}^5P_n\in\Z$, the
$2$-adic equality $\ord_2\den(P_n)=2+5\,\ord_2\dn{n}$, and the
$\hat P$ companion --- are verified exactly for $n\le8$.

Moreover the $\{2,3\}$-excess enters \emph{only} through the Laurent
coefficients $1,\tfrac16,\tfrac1{90},\dots$ of the sine kernel: the
constant term $P_n$ picks up exactly one doubled factor of the
$\zeta(2)$-coefficient $\tfrac16$, i.e.\ the factor
$12=2\cdot 6$---the concrete avatar of Lemma~\ref{lem:cost}---while
$\tfrac1{90}$ (which would inject a $5$) provably does not contribute,
matching Theorem~\ref{thm:support} and the data. The remaining gap to a
full proof of $12\,\dn{n}^5P_n\in\Z$ is a classical-type
$\dn{n}^5$-collapse bookkeeping for the structured sums above; it is
verified exactly for $n\le5$.

\section{Cross-family controls}\label{sec:crossfamily}

The same measurement standard applied to the classical corpus shows its
correction factors are uniformly \emph{removable}, in contrast to
\S\ref{sec:conjecture}:
(i) Ap\'ery $\zeta(3)$: here removability is a \emph{theorem}, not
merely our finite check: Rhin--Viola \cite[Thm.~2.1]{RV3} prove
$\dn{n}^3a_n\in\Z$ exactly (the genuine $2$ rides on the
$\zeta(3)$-coefficient parity, not the denominator); our $n\le50$
verification is a sanity check of their statement.
(ii) $\zeta(4)$: $\dn{n}^4v_n\in\Z$ for $n\le5$ in Zudilin's derived
series (the factor $2$ retained in \cite[Th\'eor\`eme 3]{KR} is slack
there).
(iii) Zudilin's asymmetric $\zeta(5),\dots,\zeta(11)$ forms: computing
the five coefficients $z_{j,n}$ exactly for $n\le5$ (apparently for the
first time), the leading factor $2$ retained even in the
\emph{conjectured} law \cite[eq.~(17.1)]{KR} is unnecessary, with
$\sim30$ bits of $2$-adic margin; the well-poised parity cancellation
and the proven inclusion of \cite{KR} serve as engine validations. For (ii) and (iii)
removability is genuinely open (Krattenthaler--Rivoal prove their
inclusions only ``\`a un facteur $2$ pr\`es'' and could not remove it);
our checks are finite verifications, and the structural reason to expect
removability there is identifiable but heuristic: the factor $2$ in the
proofs arises from termwise $2$-adic over-counting at half-integer
shifts, which the full sums have no reason to respect. What the controls
establish is a contrast, not a theorem: in every classical family
examined the correction is unneeded in all computed cases, while in the
Brown--Zudilin family---the only construction here that performs a
$\zeta(2)$-elimination---it is provably needed (the printed values of
\S\ref{sec:corrections} already require it).

\section{The worthiness supremum, and what the arithmetic cannot do}
\label{sec:gamma}

For completeness we record the analytic side. Implementing the
worthiness exponent $\gamma(\mathbf a)$ of \cite{BZ} (validated to $8$
decimals against all three published values, including every printed
intermediate), a global search of the $7$-dimensional projective cone
---steepest ascent from the published point on integer and half-integer
lattices, $120$ random multi-starts, $60$ basin hops, and an exhaustive
scan of all $31{,}710$ projectively distinct arithmetic-progression
points to height $160$---finds the supremum
\[
\sup_{\mathbf a}\gamma \;=\; 0.86597135\ldots,
\]
attained exactly at the published point $\mathbf a=(8,16,10,15,12,16,18,13)$,
whose symmetric coordinates $(20.5;\,3.5,4.5,\dots,9.5)$ form an
arithmetic progression. Combined with the bounded-slack measurements of
\S\ref{sec:conjecture}, the construction's arithmetic is now understood
to bounded factors, and none of it moves the exponential budget: within
this family the barrier to $\gamma>1$ is structural.

\section{Limitations and outlook}

All sharpness and slack statements are finite verifications under the
certification protocol of \S\ref{sec:method}; Conjecture~\ref{conj:main}
is falsifiable by a single future cell. The mechanism of
\S\ref{sec:mechanism} is proved only in the elementary
Lemma~\ref{lem:cost} plus the conditional reduction; hypothesis (H2) is
one finite lattice computation (torsion-freeness of
$H^\ast(\overline{\mathcal M}_{0,n},\Z)$ \cite{Keel} should make it
tractable), and (H1) is the integral form of the Barnes bookkeeping
sketched in \cite{BZ}. The measurement pipeline itself is portable, and
we suggest its use as an \emph{arithmetic screen} for any future
candidate constructions: the true denominators of a proposed family can
now be measured in hours rather than conjectured. In a companion
development, Lemma~4 of \cite{Zu18} (the asymptotic core of the
elementary proof that one of $\zeta(5),\dots,\zeta(25)$ is irrational)
has been fully formalized in Lean~4/Mathlib, with a complete
formalization of the theorem in progress; details will appear
separately.

\begin{thebibliography}{99}
\bibitem{BZ} F.~Brown, W.~Zudilin, \emph{On cellular rational
approximations to $\zeta(5)$}, arXiv:2210.03391.
\bibitem{Zu18} W.~Zudilin, \emph{One of the odd zeta values from
$\zeta(5)$ to $\zeta(25)$ is irrational. By elementary means}, SIGMA 14
(2018), 028; arXiv:1801.09895.
\bibitem{KR} C.~Krattenthaler, T.~Rivoal, \emph{Hyperg\'eom\'etrie et
fonction z\^eta de Riemann}, Mem.\ Amer.\ Math.\ Soc.\ 186 (2007), no.\
875; arXiv:math/0311114.
\bibitem{RV3} G.~Rhin, C.~Viola, \emph{The group structure for
$\zeta(3)$}, Acta Arith.\ 97 (2001), 269--293.
\bibitem{BrownMTZ} F.~Brown, \emph{Mixed Tate motives over $\Z$}, Ann.\
of Math.\ 175 (2012), 949--976; arXiv:1102.1312.
\bibitem{BrownP1} F.~Brown, \emph{Motivic periods and
$\mathbb P^1\setminus\{0,1,\infty\}$}, arXiv:1407.5165.
\bibitem{BrownM08} F.~Brown, \emph{Multiple zeta values and periods of
moduli spaces $\overline{\mathcal M}_{0,n}$}, Ann.\ Sci.\ \'Ec.\ Norm.\
Sup\'er.\ 42 (2009), 371--489; arXiv:math/0606419.
\bibitem{DG} P.~Deligne, A.~Goncharov, \emph{Groupes fondamentaux
motiviques de Tate mixte}, Ann.\ Sci.\ \'Ec.\ Norm.\ Sup\'er.\ 38
(2005), 1--56.
\bibitem{Schnetz} O.~Schnetz, \emph{Multiple Deligne values: a data mine
with empirically tamed denominators}, arXiv:1409.7204.
\bibitem{LS} R.~Lee, R.~Szczarba, \emph{The group $K_3(\Z)$ is cyclic of
order forty-eight}, Ann.\ of Math.\ 104 (1976), 31--60.
\bibitem{Keel} S.~Keel, \emph{Intersection theory of moduli space of
stable $n$-pointed curves of genus zero}, Trans.\ Amer.\ Math.\ Soc.\
330 (1992), 545--574.
\end{thebibliography}

\end{document}
